\documentclass[12pt, a4paper]{article}

% --- Pacotes ---
\usepackage[utf8]{inputenc}
\usepackage[T1]{fontenc}
\usepackage[brazilian]{babel}
\usepackage{graphicx}
\usepackage{float}
\usepackage{booktabs}
\usepackage{amsmath}
\usepackage{hyperref}
\usepackage[margin=2.5cm]{geometry}
\usepackage{caption}
\usepackage{subcaption}

% --- Caminho das imagens ---
\graphicspath{{imagens/}}

% --- Dados do documento ---
\title{\textbf{Lab. 2 -- Cria\c{c}\~ao de Classificadores baseado em Dados} \\[0.5cm]
\large CMC-13 -- Introdu\c{c}\~ao a Ci\^encia de Dados}
\author{
  Nome do Membro 1 \\
  Nome do Membro 2 \\
  Nome do Membro 3 \\
  Nome do Membro 4
}
\date{\today}

% =============================================================================
\begin{document}
\maketitle

% -----------------------------------------------------------------------------
\section{Prepara\c{c}\~ao dos Dados}

Descreva aqui os procedimentos drealizados para preparar os dados:
an\'alise explorat\'oria, tratamento de dados faltantes, remo\c{c}\~ao de colunas
irrelevantes, codifica\c{c}\~ao ded varssi\'aveis categ\'oricas, normaliza\c{c}\~ao, etc.

% Exemplo de como incluir figura:
% \begin{figure}[H]
%   \centering
%   \includegraphics[width=0.8\textwidth]{1_prep_correlacao_pearson.png}
%   \caption{Matriz de correla\c{c}\~ao de Pearson.}
%   \label{fig:correlacao}
% \end{figure}

% -----------------------------------------------------------------------------
\section{Cria\c{c}\~ao dos Modelos de Classifica\c{c}\~ao}

\subsection{Modelo Tradicional (KNN / \'Arvore de Decis\~ao / SVM)}

Descreva o modelo escolhsssido, hiperpar\^ametros utilizados e resultados obtidos.

\subsection{Modelo baseado em Redes Neurais (MLP)}
sss
Descreva a arquitetura da rede, fun\c{c}\~oes de ativa\c{c}\~ao, n\'umero de camadas/neur\^onios
e resultados obtidos.

\subsection{Modelo baseado em Comit\^es (Random Forest / AdaBoost / XGBoost)}

Descreva o modelo de comit\^e esssscolhido, hiperpar\^ametros e resultados obtidos.
sss
% -----------------------------------------------------------------------------
\section{An\'alise Comparativa do Desempenho dos Modelos}
ss
Apresente tabelas e gr\'aficos comparando os modelos. Utilize as m\'etricas:
acur\'acia, precision, recall, F1-score e Kappa statistic.

Discuta os resultados nos dados de treinamento e teste. H\'a varia\c{c}\~ao
significativa? Justifique.

% Exemplo de tabela comparativa:
% \begin{table}[H]
%   \centering
%   \caption{Compara\c{c}\~ao de desempenho dos modelos no conjunto de teste.}
%   \begin{tabular}{lccccc}
%     \topruleddddddddddddddddddddddddddddd
%     Modelo & Acur\'acia & Precision & Recall & F1-score & Kappa \\
%     \midrule
%     KNN         & 0.00 & 0.00 & 0.00 & 0.00 & 0.00 \\
%     MLP         & 0.00 & 0.00 & 0.00 & 0.00 & 0.00 \\
%     Random Forest & 0.00 & 0.00 & 0.00 & 0.00 & 0.00 \\
%     \bottomrule
%   \end{tabular}
%   \label{tab:comparacao}
% \end{table}

% -----------------------------------------------------------------------------
\section{Conclus\~oes, Coment\'arios e Sugest\~oes sobre o Trabalho}

Apresente as conclus\~oes do trabalho, qual modelo se mostrou mais apropriado
e sugest\~oes de melhorias.

\end{document}
